% =====================================================================
%  Lista de siglas e abreviaturas
% =====================================================================
\begin{orientacao}
Só entram aqui as siglas que aparecem no texto. Apague as que não usar e
acrescente as do seu projeto. Na primeira ocorrência no corpo do
relatório, escreva por extenso seguida da sigla entre parênteses.
Na AP2 entraram as siglas de comunicação usadas no Capítulo~5; confira,
ao fechar o relatório, se todas continuam aparecendo no texto.
\end{orientacao}

\begin{siglas}
  \item[ACK] \emph{Acknowledgement} --- confirmação de recebimento
  \item[ADC] \emph{Analog-to-Digital Converter} --- conversor analógico-digital
  \item[AIoT] \emph{Artificial Intelligence of Things}
  \item[ANVISA] Agência Nacional de Vigilância Sanitária
  \item[AP] Avaliação parcial (AP1, AP2, AP3)
  \item[API] \emph{Application Programming Interface}
  \item[BLE] \emph{Bluetooth Low Energy}
  \item[CPU] \emph{Central Processing Unit} --- processador
  \item[CSV] \emph{Comma-Separated Values}
  \item[CT] Caso de teste
  \item[DHCP] \emph{Dynamic Host Configuration Protocol}
  \item[DNS] \emph{Domain Name System}
  \item[DP] Decisão de projeto
  \item[EMA] \emph{Exponential Moving Average} --- média móvel exponencial
  \item[GATT] \emph{Generic Attribute Profile}
  \item[GPIO] \emph{General Purpose Input/Output}
  \item[HSP] Horas de sol pleno
  \item[HTTP] \emph{Hypertext Transfer Protocol}
  \item[HTTPS] HTTP sobre TLS
  \item[I\textsuperscript{2}C] \emph{Inter-Integrated Circuit}
  \item[IA] Inteligência artificial
  \item[ID] Identificador
  \item[IDE] \emph{Integrated Development Environment} --- ambiente de desenvolvimento
  \item[IEEE] \emph{Institute of Electrical and Electronics Engineers}
  \item[IoT] \emph{Internet of Things} --- internet das coisas
  \item[IP] \emph{Internet Protocol}
  \item[JSON] \emph{JavaScript Object Notation}
  \item[LED] \emph{Light-Emitting Diode} --- diodo emissor de luz
  \item[LoRaWAN] \emph{Long Range Wide Area Network}
  \item[LTE-M] \emph{Long-Term Evolution for Machines}
  \item[MQTT] \emph{Message Queuing Telemetry Transport}
  \item[MTU] \emph{Maximum Transmission Unit} --- maior unidade de dados de um pacote
  \item[NaN] \emph{Not a Number} --- valor inválido de ponto flutuante
  \item[NB-IoT] \emph{Narrowband Internet of Things}
  \item[NTP] \emph{Network Time Protocol}
  \item[NVS] \emph{Non-Volatile Storage} --- memória não volátil do ESP32
  \item[PDF] \emph{Portable Document Format}
  \item[QoS] \emph{Quality of Service} --- nível de garantia de entrega do MQTT
  \item[RAM] \emph{Random Access Memory} --- memória de trabalho, perdida ao desligar
  \item[RDC] Resolução da Diretoria Colegiada (ANVISA)
  \item[RF] Requisito funcional
  \item[RNF] Requisito não funcional
  \item[RSSI] \emph{Received Signal Strength Indicator}
  \item[RTC] \emph{Real-Time Clock} --- domínio de memória que sobrevive ao \emph{deep sleep}
  \item[SF] \emph{Spreading Factor} --- fator de espalhamento do LoRa
  \item[TCP] \emph{Transmission Control Protocol}
  \item[TI] Tecnologia da informação
  \item[TLS] \emph{Transport Layer Security}
  \item[TTN] \emph{The Things Network}
  \item[UART] \emph{Universal Asynchronous Receiver/Transmitter}
  \item[UR] Umidade relativa
  \item[USB] \emph{Universal Serial Bus}
  \item[UTC] \emph{Coordinated Universal Time} --- tempo universal coordenado
  \item[UUID] \emph{Universally Unique Identifier}
\end{siglas}
